\heading{Seminar 3}{Two routes through the consumer-theory square}

\paragraph{Conventions.} Quantities are nonnegative and $p_1,p_2,m>0$.
Targets satisfy $\bar u>0$ in problem 1 and $\bar u>1$ in problem 2 unless
stated otherwise. Write $X,H$ for demand correspondences and $x,h$ when
demand is unique. Use derivative identities where their derivatives exist.

\section*{Warm-up: true or false}
There are exactly two goods. Preferences are continuous and locally
nonsatiated; ordinary demand is unique, interior and differentiable near
the point considered. Use the compensated-demand identities where defined.
All classifications below are local, at the same prices and income:
normal means $x_{i,m}>0$, inferior means $x_{i,m}<0$, and Giffen means
$x_{i,p_i}>0$. Explain each verdict briefly.
\begin{enumerate}[label=\arabic*.]
\item Both goods can be Giffen.
\item Both goods can be normal.
\item Both goods can be inferior.
\item With $p_2=1$, a consumer's ordinary demand is locally
\[
 x_1(p_1,m)=\frac{12-m}{2-p_1},\qquad
 1<p_1<\frac43,\quad 6p_1<m<12.
\]
Your friend says: ``Demand rises with its own price, so this cannot come
from utility maximisation: it violates the law of compensated demand.''
Your friend is right.
\item After a small increase in $p_1$, with $p_2$ and income unchanged,
the consumer buys less of good 2. Therefore good 2 is inferior.
\end{enumerate}

\clearpage
\section*{1. Two papers, one research budget}
A professor is revising two papers. Let $x_i$ be hours of focused work on
paper $i$. Each hour uses $p_i$ units of a research budget $m$ (research
assistance, computing and data). The papers' quality scores are
$Q_1=\sqrt{x_1}$ and $Q_2=\sqrt{x_2}$: more work helps, but by less and
less. The second paper counts twice towards the professor's research
objective, so
\[
 u(x_1,x_2)=\sqrt{x_1}+2\sqrt{x_2}.
\]
\begin{enumerate}[label=(\alph*)]
\item Establish existence, uniqueness and budget exhaustion. Why must both
optimal quantities be strictly positive?
\item Solve utility maximisation directly: find $x_1(p,m)$, $x_2(p,m)$ and
$v(p,m)$ for general prices and income.
\item Solve expenditure minimisation directly: find $h_1(p,\bar u)$,
$h_2(p,\bar u)$ and $e(p,\bar u)$. Try $z_i=\sqrt{x_i}$.
\item Recover $x$ using Roy and $h$ using Shephard. Verify
\[
 e(p,v(p,m))=m,\qquad v(p,e(p,\bar u))=\bar u,
 \qquad h(p,v(p,m))=x(p,m).
\]
\item At $p=(1,4)$ and $m=8$, find $x$ and $v$. Calculate
$x_{1,p_1}$, $x_{1,m}$ and $h_{1,p_1}$ at the matching \emph{fixed} target.
Check the own-price Slutsky equation numerically.
\item Compute the full Slutsky matrix from Marshallian demand:
\[
 S_{ij}=\frac{\partial x_i}{\partial p_j}
       +x_j\frac{\partial x_i}{\partial m}.
\]
Check it against $D_ph$ at the matching target. Verify symmetry,
$Sp=0$ and negative semidefiniteness. Give its value at the point in (e).
\item Generalise. For an $n$-good consumer with continuous LNS utility,
positive prices and an attainable target above $u(0)$, suppose demands are
single-valued and differentiable and $e$ is $C^2$ in prices. At matching
parameters, prove that $S$ is symmetric, $Sp=0$ and $z^\top Sz\le0$.
Use Shephard and homogeneity; derive price concavity of $e$ from its
minimisation definition.
\end{enumerate}

\clearpage
\section*{2. Glitch: coffee, cake, and a head start}
At Glitch, the NES caf\'e, a student buys coffee ($x_1$) and cake ($x_2$),
measured in divisible serving-equivalents, at prices $p_1$ and $p_2$.
The student has a break budget $m$. A familiar coffee ritual has a one-point
head start; a cake break has diminishing returns. For this student, what
matters is whichever of the two routes provides the better break:
\[
 u(x_1,x_2)=\max\{x_1+1,\sqrt{x_2}\}.
\]
\begin{enumerate}[label=(\alph*)]
\item Check continuity, local nonsatiation and convexity of preferences.
Draw $\{x\ge0:u(x)\ge2\}$. What is $u(0)$, and why do we work with targets
$\bar u>1$ for the main duality calculations?
\item Find $v(p,m)$ and the entire $X(p,m)$ by comparing the two corner
routes. At $p=(9/2,1)$, find every income range in which each route is optimal,
including both switching boundaries.
\item Find $e(p,\bar u)$ and the entire $H(p,\bar u)$ for $\bar u>1$.
At $p=(9/2,1)$, find every target range and both switching targets.
What happens for targets $\bar u\le1$?
\item Apply Roy and Shephard in each strict regime. Explain what prevents
a single gradient from selecting a bundle at a switching boundary.
\item Verify both inverse-value identities and demand-set equality at matching
parameters. At the specified prices, match the two income boundaries to
the two target boundaries.
\item At $p=(9/2,1)$ and $m=9/4$, list all solutions to both corresponding
problems. Is the midpoint also optimal? Check the other boundary, $m=9$.
\item \emph{Extension.} Set $p=(4,1)$ and $m=4$. Does the derivative $v_m$
exist at this tie? Does that alone let you use Roy's identity?
\end{enumerate}
